% !TeX spellcheck = en_US
\documentclass[french]{yLectureNote}

\title{Algèbre linéaire 2}
\subtitle{Langage mathématique}
\author{Paulhenry Saux}
\date{\today}
\yLanguage{Français}

\professor{C.Dartyge}

\usepackage{graphicx}%----pour mettre des images
\usepackage[utf8]{inputenc}%---encodage
\usepackage{geometry}%---pour modifier les tailles et mettre a4paper
%\usepackage{awesomebox}%---pour les boites d'exercices, de pbq et de croquis ---d\'esactiv\'e pour les TP de PC
\usepackage{tikz}%---pour deiffner + d\'ependance de chemfig
\usepackage{tkz-tab}
\usepackage{chemfig}%---pour deiffner formules chimiques
\usepackage{chemformula}%---pour les formules chimiques en \'equation : \ch{...}
\usepackage{tabularx}%---pour dimensionner automatiquement les tableaux avec variable X
\usepackage{awesomebox}%---Pour les boites info, danger et autres
\usepackage{menukeys}%---Pour deiffner les touches de Calculatrice
\usepackage{fancyhdr}%---pour les en-t\^ete personnalis\'ees
\usepackage{blindtext}%---pour les liens
\usepackage{hyperref}%---pour les liens (\`a mettre en dernier)
\usepackage{caption}%---pour la francisation de la l\'egende table vers Tableau
\usepackage{pifont}
\usepackage{array}%---pour les tableaux
\usepackage{lipsum}
\usepackage{yFlatTable}
\usepackage{multicol}
\newcommand{\Lim}[1]{\lim\limits_{\substack{#1}}\:}
\renewcommand{\vec}{\overrightarrow}
\newcommand{\bz}{\overline{z}}
\DeclareMathOperator{\card}{card}
\renewcommand{\vec}{\overrightarrow}
\newcommand{\N}[0]{\mathbb{N}}
\newcommand{\R}[0]{\mathbb{R}}
\newcommand{\C}[0]{\mathbb{C}}
\newcommand{\dd}[0]{\mathrm{d}}
\begin{document}
\setcounter{chapter}{3}
\chapter{Diagonalisation et trigonalisation}

\section{Matrice d'Exemple $A$}

Soit la matrice $A$ définie par :
$$A = \begin{pmatrix} 2 & 1 & 1 \\ 0 & 3 & 1 \\ 0 & 0 & 2 \end{pmatrix}$$

\section{Diagonalisation d'une Matrice}

La matrice $A$ est diagonalisable si et seulement si $\sum_{\lambda} \dim(E_\lambda) = \dim(A)$. L'objectif est de trouver $P$ inversible et $D$ diagonale telles que $A = P D P^{-1}$.

\subsection{Polynôme Caractéristique $P_A(\lambda)$ et les Valeurs Propres}
\explanation{1}{Souvent on trouvera un polynome dont il faudra trouver les racines pour l'écrire sous sa forme factorisée.}
\begin{enumerate}
    \item \textbf{Calcul de $A - \lambda I$ :}
    $$A - \lambda I = \begin{pmatrix} 2 - \lambda & 1 & 1 \\ 0 & 3 - \lambda & 1 \\ 0 & 0 & 2 - \lambda \end{pmatrix}$$

    \item \textbf{Calcul du Déterminant :} Puisque $A - \lambda I$ est triangulaire\explain{1}{right}{0}{0.5}{} :
    $$P_A(\lambda) = \det(A - \lambda I) = (2 - \lambda)(3 - \lambda)(2 - \lambda) = (2 - \lambda)^2 (3 - \lambda)$$

    \item \textbf{Valeurs Propres :} Les racines de $P_A(\lambda) = 0$ sont :
    \begin{itemize}
        \item $\lambda_1 = 2$ (Multiplicité Algébrique $m_a(\lambda_1) = 2$)
        \item $\lambda_2 = 3$ (Multiplicité Algébrique $m_a(\lambda_2) = 1$)
    \end{itemize}
\end{enumerate}

\subsection{Multiplicité Géométrique $m_g(\lambda)$}

On doit vérifier si $m_g(\lambda) = m_a(\lambda)$ pour chaque $\lambda$. $m_g(\lambda) = \dim(E_\lambda) = n - \text{rang}(A - \lambda I)$.

\begin{itemize}
    \item \textbf{Cas $\lambda_2 = 3$ :} $m_a(3) = 1$. Nécessairement, $m_g(3) = 1$.

    \item \textbf{Cas $\lambda_1 = 2$ :} $m_a(2) = 2$.
    \begin{itemize}
        \item Calcul de $A - 2I$ :
        $$A - 2I = \begin{pmatrix} 0 & 1 & 1 \\ 0 & 1 & 1 \\ 0 & 0 & 0 \end{pmatrix}$$
        \item Calcul du rang : $\text{rang}(A - 2I) = 1$ (les deux premières lignes sont identiques et non nulles).
        \item Multiplicité Géométrique :
        $$m_g(2) = 3 - \text{rang}(A - 2I) = 3 - 1 = 2$$
    \end{itemize}
\end{itemize}

\textbf{Conclusion :} Puisque $m_g(2) = m_a(2) = 2$ et $m_g(3) = m_a(3) = 1$, la matrice $A$ est \textbf{diagonalisable}.

\subsection{Calculer les Vecteurs Propres}
\explanation{2}{Comme la multiplicité de \(\lambda_1\) est 2, on déduit 2 vecteurs.}
\begin{itemize}
    \item \textbf{Pour $\lambda_2 = 3$ :} $(A - 3I)v = 0$
    $$\begin{pmatrix} -1 & 1 & 1 \\ 0 & 0 & 1 \\ 0 & 0 & -1 \end{pmatrix} \begin{pmatrix} x \\ y \\ z \end{pmatrix} = \begin{pmatrix} 0 \\ 0 \\ 0 \end{pmatrix} \implies \begin{cases} -x + y + z = 0 \\ z = 0 \end{cases}$$
    $z=0$ et $-x+y=0 \implies x=y$.
    $E_3 = \text{Vect}\left( v_1 = \begin{pmatrix} 1 \\ 1 \\ 0 \end{pmatrix} \right)$.

    \item \textbf{Pour $\lambda_1 = 2$ :} $(A - 2I)v = 0$
    $$\begin{pmatrix} 0 & 1 & 1 \\ 0 & 1 & 1 \\ 0 & 0 & 0 \end{pmatrix} \begin{pmatrix} x \\ y \\ z \end{pmatrix} = \begin{pmatrix} 0 \\ 0 \\ 0 \end{pmatrix} \implies y + z = 0 \implies z = -y$$
    $x$ est libre. $E_2 = \left\{ \begin{pmatrix} x \\ y \\ -y \end{pmatrix} \right\} = \text{Vect}\left( v_2 = \begin{pmatrix} 1 \\ 0 \\ 0 \end{pmatrix}, v_3 = \begin{pmatrix} 0 \\ 1 \\ -1 \end{pmatrix} \right)$\explain{2}{right}{0}{0.5}{}.
\end{itemize}

\subsection{Former $P$ et $D$}

\begin{itemize}
    \item \textbf{Matrice de Passage $P$} (colonnes : $v_1, v_2, v_3$) :
    $$P = \begin{pmatrix} 1 & 1 & 0 \\ 1 & 0 & 1 \\ 0 & 0 & -1 \end{pmatrix}$$

    \item \textbf{Matrice Diagonale $D$} (diagonale : $\lambda_2, \lambda_1, \lambda_1$) :
    $$D = \begin{pmatrix} 3 & 0 & 0 \\ 0 & 2 & 0 \\ 0 & 0 & 2 \end{pmatrix}$$
\end{itemize}

\textbf{Résultat :} $A = P D P^{-1}$.

\end{document}


